\emph{Setting and inputs.} The unit and the loop conventions are those of the one-loop result in the spacetime algebra: the Dirac matrices are the vectors, the units of the Feynman rules are coefficients outside the Clifford product, the trace is $4\langle X\rangle_{0}$. Dimensional regularisation and its chirality element are imported (Fact \claimref{synthesis.imported-dimensional-regularisation-and-the-bmhv-gamma5}): in $\Cl(1, D-1)$ the four-dimensional $I_{4}$ anticommutes with $\gamma_{0}, \ldots, \gamma_{3}$ and commutes with the evanescent generators (Theorem \claimref{synthesis.unit-i4-keeps-its-three-properties-only-in-dimension-zero-mod-four}(iv)), so $\gamma_{5} = iI_{4}$, $i$ the coefficient, is the scheme's element: $\gamma_{5}^{2} = 1$, $\{\gamma_{5}, \bar\gamma^{\mu}\} = 0$, $[\gamma_{5}, \hat\gamma^{\mu}] = 0$. With the trace normalised by $\mathrm{Tr}\,1 = 4$ in every dimension, $\langle X\rangle_{0}$ is the scheme's normalised trace. Imported as well: Feynman parametrisation (Fact \claimref{synthesis.imported-one-loop-integrals-and-pauli-villars-regulators}(i)), and, from the dimensional-regularisation Fact (\claimref{synthesis.imported-dimensional-regularisation-and-the-bmhv-gamma5}(iv), the radial integral and symmetric integration, here after the Wick rotation; in Minkowski form Peskin and Schroeder, eqs. (A.44) to (A.47)), the symmetric integration $\ell^{a}\ell^{b} \to g^{ab}Q(\ell)/D$, $\ell^{a}\ell^{b}\ell^{c}\ell^{d} \to S^{abcd}Q(\ell)^{2}/(D(D+2))$, and the master integrals
\[
\int\frac{d^{D}\ell}{(2\pi)^{D}}\frac{Q(\ell)}{(Q(\ell) - \Delta)^{2}} = \frac{-i}{(4\pi)^{D/2}}\frac{D}{2}\,\Gamma(1 - \tfrac{D}{2})\,\Delta^{D/2 - 1}, \quad \int\frac{d^{D}\ell}{(2\pi)^{D}}\frac{Q(\ell)^{2}}{(Q(\ell) - \Delta)^{4}} = \frac{i}{(4\pi)^{D/2}}\frac{D(D+2)}{4}\frac{\Gamma(2 - \tfrac{D}{2})}{\Gamma(4)}\Delta^{D/2 - 2}.
\]
Every Clifford statement below is exact and is checked over the Gaussian rationals.

\emph{(i)} For a vector $k = \bar k + \hat k$, $P_{R}\bar k = \bar kP_{L}$ and $P_{R}\hat k = \hat kP_{R}$, and $P_{L}\bar\gamma^{\nu}P_{R} = \bar\gamma^{\nu}P_{R}$, $P_{R}\bar\gamma^{\nu}P_{R} = \bar\gamma^{\nu}P_{L}P_{R} = 0$. Hence $\bar\gamma^{\mu}P_{R}k\bar\gamma^{\nu}P_{R} = \bar\gamma^{\mu}\bar k\bar\gamma^{\nu}P_{R}$; the same with $P_{L}$; around a loop of chiral vertices the projectors pass through every numerator and leave its barred part (checked in $\Cl(1,5)$ for both chiralities, for four index pairs and for a four-vertex loop at rational momenta with evanescent components; the control without projectors differs). A remaining evanescent factor $\hat k$ next to barred vectors and $P_{R}$ has zero scalar part (checked), since the metric pairs no barred with hatted direction. In the language of the unit result, $\hat k$ commutes with $I_{4}$, so on the mixed blades $\bar\gamma^{\mu}\hat k$ the left and right units differ (its clause (v)), and the projector built from the left element annihilates them.

\emph{(ii)} Two-point function. The loop is $i\Pi^{\mu\nu} = (-1)(ieY)^{2}i^{2}\int\mathrm{Tr}[\bar\gamma^{\mu}P_{R}k\bar\gamma^{\nu}P_{R}k']/(Q(k)Q(k'))$, prefactor $-e^{2}Y^{2}$ (checked). By (i) the numerator is $\mathrm{Tr}[\bar\gamma^{\mu}\bar k\bar\gamma^{\nu}\bar k'P_{R}]$; its $\gamma_{5}$ part is an oriented volume of $\bar\gamma^{\mu}, \bar\gamma^{\nu}$ and the momenta, and with one external momentum $q$ it integrates to zero. With $k = \ell - xq$, $k' = \ell + (1-x)q$, the terms without $\ell$ coincide with those of the vector-like loop (the external momenta are barred); the terms linear in $\ell$ integrate to zero; the $\ell\ell$ term is, after symmetric integration, $\tfrac{1}{D}Q(\ell)\sum_{a}$ over barred $a$ of $\tfrac12\langle\bar\gamma^{\mu}\gamma_{a}\bar\gamma^{\nu}\gamma^{a}\rangle_{0}$ for the chiral loop and over all $a$ for half the Dirac loop. The check evaluates both at $D = 4, \ldots, 8$: $-\bar g^{\mu\nu}$ and $\tfrac{2-D}{2}\bar g^{\mu\nu}$ (degree one in $D$, so two values decide it and the others confirm). The half Dirac loop in dimensional regularisation is transverse (vector-like, Fact \claimref{synthesis.imported-dimensional-regularisation-and-the-bmhv-gamma5}(i)); the chiral loop exceeds it by $\tfrac{D-4}{2}\bar g^{\mu\nu}$ (normalised), that is, by $4\cdot\tfrac{D-4}{2}\cdot\tfrac{1}{D}Q(\ell)\bar g^{\mu\nu} = \tfrac{2(D-4)}{D}Q(\ell)\bar g^{\mu\nu}$ in the integrand of the first master integral. With its factor $\tfrac{D}{2}\Gamma(1 - \tfrac{D}{2})$ the product is $(D-4)\Gamma(1 - \tfrac{D}{2})$, and $(D-4)\Gamma(1 - \tfrac{D}{2}) = \tfrac{-2}{1 - D/2}\Gamma(3 - \tfrac{D}{2}) \to 2$ at $D = 4$ (the rational factor checked exactly, $\Gamma(1) = 1$). Hence the difference of the integrated numerators is $-2i\Delta\bar g^{\mu\nu}/(16\pi^{2})$ with $\Delta = -x(1-x)Q(q)$, and $\int_{0}^{1}x(1-x)dx = \tfrac16$; so $i\Pi_{\mathrm{break}} = -e^{2}Y^{2}\cdot(-2i)(-\tfrac16)Q(q)\bar g^{\mu\nu}/(16\pi^{2})$, $\Pi_{\mathrm{break}} = -\tfrac{e^{2}Y^{2}}{48\pi^{2}}Q(\bar q)\bar g^{\mu\nu}$ (assembled in the check from the computed contraction). It is not transverse: $q_{\mu}\Pi_{\mathrm{break}}^{\mu\nu} \propto Q(q)q^{\nu} \neq 0$. The other contributions to the difference carry a factor $D - 4$ times a convergent integral and vanish at $D = 4$.

\emph{(iii)} Four-point function. By (i) every numerator around the four-vertex loop is barred, so the $\ell^{4}$ term of the chiral loop is a four-dimensional contraction; summed over the six orderings of the photons (the $(4-1)!$ diagrams of a fermion loop) it is zero, and so is its $\gamma_{5}$ part (checked for seven index tuples; it vanishes in every $D$ since it uses barred indices only). For half the Dirac loop the same sum is $(D-4)(D-2)S^{\mu\nu\rho\sigma}$ (checked at $D = 4, 5, 6, 7$: the quotient by $(D-4)S$ is $D - 2$, linear in $D$, confirmed at three points). The vector-like loop is gauge invariant, and the remaining terms of the difference carry $D - 4$ times convergent integrals; so the breaking is $-(D-4)(D-2)S$ (normalised) times $4/(D(D+2))$ times the second master integral, $(D-4)\Gamma(2 - \tfrac{D}{2}) = -2\Gamma(3 - \tfrac{D}{2}) \to -2$ (checked), $D - 2 \to 2$, and Feynman parametrisation with $3!$ times the simplex volume $\tfrac16$ gives $1$ for a $\Delta$-independent limit. The result is $i\cdot\tfrac{2}{3}S/(16\pi^{2})$, and with the prefactor $(-1)(ieY)^{4}i^{4} = -e^{4}Y^{4}$, $\Gamma_{\mathrm{break}} = -\tfrac{2}{3}\tfrac{e^{4}Y^{4}}{16\pi^{2}}S = -\tfrac{e^{4}Y^{4}}{24\pi^{2}}S$.

\emph{(iv)} Counterterms. The effective action contains $\tfrac12\int\bar A_{\mu}\Pi^{\mu\nu}\bar A_{\nu}$ with $Q(\bar q) \to -\bar\partial^{2}$; the breaking part is $\tfrac{1}{16\pi^{2}}\tfrac{e^{2}Y^{2}}{6}\int\bar A_{\mu}\bar\partial^{2}\bar A^{\mu}$, and the counterterm is its negative. A term $c(\bar A^{2})^{2} = \tfrac{c}{3}S_{\mu\nu\rho\sigma}\bar A^{\mu}\bar A^{\nu}\bar A^{\rho}\bar A^{\sigma}$ has four-point vertex $8cS$, so the counterterm cancelling $\Gamma_{\mathrm{break}}$ has $c = \tfrac{e^{4}Y^{4}}{12}\cdot\tfrac{1}{16\pi^{2}}$. Each field contributes separately, so the sums over fields are $\sum Y^{2}$ and $\sum Y^{4}$. The coefficients $-\tfrac16$ and $\tfrac{1}{12}$ (in units $e^{2}\mathrm{Tr}\,\mathcal{Y}_{R}^{2}/(16\pi^{2})$ and $e^{4}\mathrm{Tr}\,\mathcal{Y}_{R}^{4}/(16\pi^{2})$) are those of \cite{BeluscaMaito2021}, eq. (6.5), whose eq. (6.4) gives the corresponding breaking $\tfrac{e^{2}\mathrm{Tr}\mathcal{Y}_{R}^{2}}{3}(\bar\partial_{\mu}c)(\bar\partial^{2}\bar A^{\mu})$, the variation of the first term under $\delta\bar A = \bar\partial c$; the agreement is a check on the whole chain, signs included. For the left-handed content of Theorem \claimref{synthesis.anomaly-freedom-forces-slot-content}(ii), $m = 1$, the colour and weak multiplicities give $\sum Y^{2} = \tfrac{6}{36} + \tfrac{12}{9} + \tfrac{3}{9} + \tfrac{2}{4} + 1 = \tfrac{10}{3}$ and $\sum Y^{4} = \tfrac{95}{54}$, and $\sum Y^{3} = 0$ (checked); the hypercharge counterterm coefficients are $-\tfrac59$ and $\tfrac{95}{648}$ in units $g'^{2}/(16\pi^{2})$ and $g'^{4}/(16\pi^{2})$ per generation. Since $Y^{2}, Y^{4} \geq 0$, in the purely chiral formulation (every field with the vertex $\bar\gamma^{\mu}P$) these never cancel between fields; a vector-like pair given a $D$-dimensional vector vertex has no such breaking. The one-loop structure is standard \cite{BeluscaMaito2020}. This is in contrast with the triangle, whose non-removable breaking in this scheme is the anomaly, proportional to $\sum Y^{3}$ (standard; \cite{BeluscaMaito2021} impose $\mathrm{Tr}\,\mathcal{Y}_{R}^{3} = 0$, so their eq. (6.4) has no such term), computed with Pauli-Villars fields in \claimref{synthesis.triangle-anomaly-in-the-spacetime-algebra-is-the-pseudoscalar-part-of-the-field-square-and-cancels-per-generation} and scheme independent by \claimref{synthesis.one-loop-vacuum-polarisation-extends-without-a-regulator-and-ward-preserving-extensions-differ-only-by-the-renormalisation-constants}(iii), and which vanishes on one generation. $\square$
